\section{Applications of Sequences}

\begin{toolkitbox}{Sequence Applications --- Quick Reference}
\small
\begin{tabular}{l l}
\toprule
\textbf{Construction} & \textbf{Role} \\
\midrule
\hyperref[def:newton-sequence-sqrt-two]{Newton sequence for \(\sqrt2\)} & Iterative approximation model \\
\hyperref[thm:newton-approximation-sqrt-two]{Newton approximation} & Converges to \(\sqrt2\) \\
\hyperref[def:factorial-partial-sums]{Factorial partial sums} & Series approximants for \(e\) \\
\hyperref[thm:factorial-partial-sums-approximate-e]{Approximation of \(e\)} & Factorial sums converge to \(e\) \\
\hyperref[def:compound-interest-sequence]{Compound-interest sequence} & Classical exponential approximants \\
\hyperref[thm:compound-interest-approximation-e]{Compound-interest approximation} & Compound interest converges to \(e\) \\
\hyperref[def:decimal-truncation-sequence]{Decimal truncation sequence} & Decimal approximants from below \\
\hyperref[thm:decimal-truncations-converge]{Decimal truncations} & Truncations converge to the target real \\
\bottomrule
\end{tabular}
\end{toolkitbox}

\subsection*{Applications of Sequences}
\begin{exposition}
The convergence theorems for sequences become numerical tools once the
sequence is computable. Newton iteration turns an equation into a convergent
recursive approximation; factorial partial sums and compound-interest terms
approximate the Euler number; decimal truncations approximate arbitrary real
numbers by finite decimals.
\end{exposition}

% ---------------------------------------------------------
% Definition --- Newton Sequence for sqrt(2)
% ---------------------------------------------------------

\begin{definitionbox}{Definition (Newton Sequence for $\sqrt{2}$)}
\begin{definition}[Newton Sequence for $\sqrt{2}$]
\label{def:newton-sequence-sqrt-two}
The \textbf{Newton sequence for} $\sqrt{2}$ is the real sequence $(x_n)$
defined by
\[
  x_1:=\frac{3}{2},
  \qquad
  x_{n+1}:=\frac{1}{2}\left(x_n+\frac{2}{x_n}\right)
  \quad\text{for every } n\in\mathbb{N}.
\]
\end{definition}
\end{definitionbox}

\begin{remark*}[Standard quantified statement]
\[
\begin{aligned}
&x_1=\frac{3}{2}\\
&\quad\land\quad
\forall n\in\mathbb{N}\;
\left(
  x_{n+1}=\frac{1}{2}\left(x_n+\frac{2}{x_n}\right)
\right).
\end{aligned}
\]
\end{remark*}

\begin{remark*}[Interpretation]
The recursion is Newton's method applied to the equation $x^2-2=0$.
Each new term averages the current estimate with the correction $2/x_n$.
\end{remark*}

\begin{dependencies}
\begin{itemize}
  \item \hyperref[def:real-sequence]{Real Sequence}
\end{itemize}
\end{dependencies}

% ---------------------------------------------------------
% Theorem --- Newton Approximation of sqrt(2)
% ---------------------------------------------------------

\begin{theorembox}{Theorem (Newton Approximation of $\sqrt{2}$)}
\begin{theorem}[Newton Approximation of $\sqrt{2}$]
\label{thm:newton-approximation-sqrt-two}
Let $(x_n)$ be the Newton sequence for $\sqrt{2}$. Then $(x_n)$ converges
and
\[
  \lim_{n\to\infty}x_n=\sqrt{2}.
\]

\smallskip
\noindent
\hyperref[prf:newton-approximation-sqrt-two]{\textit{Go to proof.}}
\end{theorem}
\end{theorembox}

\begin{remark*}[Standard quantified statement]
\[
\forall \varepsilon>0\;\exists N\in\mathbb{N}\;\forall n\ge N
\;\bigl(|x_n-\sqrt{2}|<\varepsilon\bigr).
\]
\end{remark*}

\begin{remark*}[Predicate reading]
\[
\operatorname{ConvergentSequence}(x_n,\sqrt{2}).
\]
\end{remark*}

\begin{remark*}[Interpretation]
The Newton sequence is decreasing after the first term and bounded below by
$\sqrt{2}$, so monotone convergence supplies a limit. Passing the recursion
to the limit identifies the limit as the positive solution of $L^2=2$.
\end{remark*}

\begin{dependencies}
\begin{itemize}
  \item \hyperref[def:newton-sequence-sqrt-two]{Newton Sequence for $\sqrt{2}$}
  \item \hyperref[thm:monotone-convergence-theorem]{Monotone Convergence Theorem}
  \item \hyperref[thm:rational-sequence-limit]{Rational Sequence Limit}
  \item \hyperref[def:square-root-sequence]{Square Root Sequence}
\end{itemize}
\end{dependencies}

% ---------------------------------------------------------
% Definition --- Factorial Partial Sums
% ---------------------------------------------------------

\begin{definitionbox}{Definition (Factorial Partial Sums)}
\begin{definition}[Factorial Partial Sums]
\label{def:factorial-partial-sums}
The \textbf{factorial partial-sum sequence} is the real sequence $(s_n)$
defined by
\[
  s_n:=\sum_{k=0}^{n}\frac{1}{k!}
  \qquad\text{for every } n\in\mathbb{N}.
\]
\end{definition}
\end{definitionbox}

\begin{remark*}[Standard quantified statement]
\[
\forall n\in\mathbb{N}\;
\left(
  s_n=\sum_{k=0}^{n}\frac{1}{k!}
\right).
\]
\end{remark*}

\begin{remark*}[Interpretation]
The sequence adds the reciprocal factorial terms one at a time. The terms
decrease very rapidly, so the partial sums stabilize quickly.
\end{remark*}

\begin{dependencies}
\begin{itemize}
  \item \hyperref[def:real-sequence]{Real Sequence}
\end{itemize}
\end{dependencies}

% ---------------------------------------------------------
% Theorem --- Factorial Partial Sums Approximate e
% ---------------------------------------------------------

\begin{theorembox}{Theorem (Factorial Partial Sums Approximate $e$)}
\begin{theorem}[Factorial Partial Sums Approximate $e$]
\label{thm:factorial-partial-sums-approximate-e}
Let $(s_n)$ be the factorial partial-sum sequence. Then $(s_n)$ converges
to the Euler number $e$:
\[
  \lim_{n\to\infty}s_n=e.
\]

\smallskip
\noindent
\hyperref[prf:factorial-partial-sums-approximate-e]{\textit{Go to proof.}}
\end{theorem}
\end{theorembox}

\begin{remark*}[Standard quantified statement]
\[
\forall \varepsilon>0\;\exists N\in\mathbb{N}\;\forall n\ge N
\;\bigl(|s_n-e|<\varepsilon\bigr).
\]
\end{remark*}

\begin{remark*}[Predicate reading]
\[
\operatorname{ConvergentSequence}(s_n,e).
\]
\end{remark*}

\begin{remark*}[Interpretation]
This is the sequence form of the factorial-series definition of $e$. In
practice the approximation is efficient because the error after $n$ terms is
controlled by a rapidly shrinking factorial tail.
\end{remark*}

\begin{dependencies}
\begin{itemize}
  \item \hyperref[def:factorial-partial-sums]{Factorial Partial Sums}
  \item \hyperref[def:number-e]{The Number $e$}
  \item \hyperref[thm:cauchy-criterion-real-sequences]{Cauchy Criterion for Real Sequences}
\end{itemize}
\end{dependencies}

% ---------------------------------------------------------
% Definition --- Compound-Interest Sequence
% ---------------------------------------------------------

\begin{definitionbox}{Definition (Compound-Interest Sequence)}
\begin{definition}[Compound-Interest Sequence]
\label{def:compound-interest-sequence}
The \textbf{compound-interest sequence} is the real sequence $(c_n)$
defined by
\[
  c_n:=\left(1+\frac{1}{n}\right)^n
  \qquad\text{for every } n\in\mathbb{N}.
\]
\end{definition}
\end{definitionbox}

\begin{remark*}[Standard quantified statement]
\[
\forall n\in\mathbb{N}\;
\left(
  c_n=\left(1+\frac{1}{n}\right)^n
\right).
\]
\end{remark*}

\begin{remark*}[Interpretation]
The term $c_n$ is the value of one unit of principal after one unit of time
when interest is compounded $n$ times at nominal rate $1$.
\end{remark*}

\begin{dependencies}
\begin{itemize}
  \item \hyperref[def:real-sequence]{Real Sequence}
\end{itemize}
\end{dependencies}

% ---------------------------------------------------------
% Theorem --- Compound-Interest Approximation of e
% ---------------------------------------------------------

\begin{theorembox}{Theorem (Compound-Interest Approximation of $e$)}
\begin{theorem}[Compound-Interest Approximation of $e$]
\label{thm:compound-interest-approximation-e}
Let $(c_n)$ be the compound-interest sequence. Then
\[
  \lim_{n\to\infty}c_n=e.
\]

\smallskip
\noindent
\hyperref[prf:compound-interest-approximation-e]{\textit{Go to proof.}}
\end{theorem}
\end{theorembox}

\begin{remark*}[Standard quantified statement]
\[
\forall \varepsilon>0\;\exists N\in\mathbb{N}\;\forall n\ge N
\;\bigl(|c_n-e|<\varepsilon\bigr).
\]
\end{remark*}

\begin{remark*}[Predicate reading]
\[
\operatorname{ConvergentSequence}(c_n,e).
\]
\end{remark*}

\begin{remark*}[Interpretation]
The compound-interest approximation gives a second natural sequence
converging to the same constant as the factorial partial sums. It is slower
numerically, but it explains why $e$ appears in continuous growth.
\end{remark*}

\begin{dependencies}
\begin{itemize}
  \item \hyperref[def:compound-interest-sequence]{Compound-Interest Sequence}
  \item \hyperref[def:number-e]{The Number $e$}
  \item \hyperref[thm:sequence-squeeze-theorem]{Sequence Squeeze Theorem}
\end{itemize}
\end{dependencies}

% ---------------------------------------------------------
% Definition --- Decimal Truncation Sequence
% ---------------------------------------------------------

\begin{definitionbox}{Definition (Decimal Truncation Sequence)}
\begin{definition}[Decimal Truncation Sequence]
\label{def:decimal-truncation-sequence}
Let $\alpha\in\mathbb{R}$. The \textbf{decimal truncation sequence} of
$\alpha$ is the real sequence $(d_n)$ defined by
\[
  d_n:=\frac{\lfloor 10^n\alpha\rfloor}{10^n}
  \qquad\text{for every } n\in\mathbb{N}.
\]
\end{definition}
\end{definitionbox}

\begin{remark*}[Standard quantified statement]
\[
\forall n\in\mathbb{N}\;
\left(
  d_n=\frac{\lfloor 10^n\alpha\rfloor}{10^n}
\right).
\]
\end{remark*}

\begin{remark*}[Interpretation]
The number $d_n$ is obtained by cutting $\alpha$ off after $n$ decimal
places. Truncation is one-sided: it does not round to the nearest decimal.
\end{remark*}

\begin{dependencies}
\begin{itemize}
  \item \hyperref[def:real-sequence]{Real Sequence}
\end{itemize}
\end{dependencies}

% ---------------------------------------------------------
% Theorem --- Decimal Truncations Converge
% ---------------------------------------------------------

\begin{theorembox}{Theorem (Decimal Truncations Converge)}
\begin{theorem}[Decimal Truncations Converge]
\label{thm:decimal-truncations-converge}
Let $\alpha\in\mathbb{R}$, and let $(d_n)$ be the decimal truncation
sequence of $\alpha$. Then
\[
  \lim_{n\to\infty}d_n=\alpha.
\]

\smallskip
\noindent
\hyperref[prf:decimal-truncations-converge]{\textit{Go to proof.}}
\end{theorem}
\end{theorembox}

\begin{remark*}[Standard quantified statement]
\[
\forall \varepsilon>0\;\exists N\in\mathbb{N}\;\forall n\ge N
\;\bigl(|d_n-\alpha|<\varepsilon\bigr).
\]
\end{remark*}

\begin{remark*}[Predicate reading]
\[
\operatorname{ConvergentSequence}(d_n,\alpha).
\]
\end{remark*}

\begin{remark*}[Interpretation]
Decimal truncations converge because the truncation error is always less
than $10^{-n}$. This is one of the most concrete examples of convergence:
finite decimal data can approximate a real number to arbitrary accuracy.
\end{remark*}

\begin{dependencies}
\begin{itemize}
  \item \hyperref[def:decimal-truncation-sequence]{Decimal Truncation Sequence}
  \item \hyperref[thm:sequence-squeeze-theorem]{Sequence Squeeze Theorem}
\end{itemize}
\end{dependencies}

